\section{Discussion}
\label{sec:discussion}

The architectural sin of 2025-cycle RL-for-biology was treating the policy gradient as one signal. PRICE-RL splits it into two algebraically identifiable components --- selection and transmission --- the same two that the Price equation imposes on evolutionary change and that \citet{frank2025price} proved are universal under the force--metric--bias law. The split is exact (Theorem~\ref{thm:exact}, verified to cosine $1.0000$ in $1{,}600$ rounds across the large NK sweep alone), the resulting Price ratio is a strictly stronger reward-hacking diagnostic than the current state of the art (mean lead $18.6$ rounds on $100\%$ of $16$ seeds, robust to surrogate noise up to $\sigma_n = 0.40$), and adaptive mixing of the components --- with a target derived from the reward landscape's autocorrelation length --- ties AdaLead at small GB1 budgets and beats it by $12.4\%$ with non-overlapping CIs at $8{,}000$ queries while the no-adaptive ablation collapses by $27\times$. Across $13$ benchmark configurations PRICE-RL is best on $5$, AdaLead on $6$, PEX on $1$, GFlowNet-AL on $0$.

The two findings that move the field forward most --- in our reading --- are the Trap-5 result and the reward-hacking diagnostic, not the GB1 numbers. Trap-5 is a binary deceptive landscape with no biological content, where AdaLead misses the global optimum on $\sim 5\%$ of seeds at $N=30$ and up to $20\%$ at $N=120$, and where PRICE-RL hits the optimum on every seed at every scale. The Price decomposition is not a biology-specific trick; it is a general RL primitive that finds its cleanest empirical demonstration in biology because biological optimisation \emph{is} a Price-equation process at the level of populations and generations, and our algorithm imposes the same algebra at the level of policies and rounds. The reward-hacking diagnostic is per-round and gradient-level, in contrast to the post-hoc behaviour-level statistics used by every prior work, and it transferred unchanged to a token-level autoregressive RLHF setup --- evidence that the same diagnostic should fire on LLM reward-model hacking strictly before any behavioural metric.

The honest negatives are equally informative. T2 (PI tight tracking) is structurally bounded: $\rho_t = \|g_S\|/(\|g_S\|+\|g_T\|)$ depends on the in-support fraction of samples, which is a property of the policy distribution rather than a free knob, and closure attempts based on entropy injection, autoregressive policy capacity, and an entropy-target redesign all fail in essentially the same way. Long-sequence DMS at scale is the second honest negative: WT-aware initialisation and a locality kernel close most but not all of the gap, and the proposal's ESM-2 backbone is the canonical solution. The closed-loop AL setting --- in which the policy trains on a surrogate retrained each round, as in $\delta$-CS, $\mu$Protein, and ALDE --- favours AdaLead because the mutate-from-best operator is invariant to surrogate noise while PRICE-RL inherits surrogate-prediction error in its gradient signal; a naïve $\sigma$-gate is the wrong fix. The $\rho^*(L)$ direction inversion is also worth flagging: our autocorrelation-to-target mapping is the \emph{opposite} sign of what a high-data regret analysis would predict, and we arrived at it by revising based on empirics ($r=-0.79$ before, $r=0.995$ after) rather than a priori derivation.

Looking forward: a non-factorised autoregressive policy class is the natural next step both for tightening $\rho_t$ tracking and for cross-protein transfer, where the factorised representation entangles selection rank-structure and transmission concentration in a single softmax. The Price ratio's status as an early-warning signal for reward hacking transfers naturally to instruction-tuning RLHF (tokens are the action space, the support test is at the per-token softmax). The multi-objective extension --- where the selection covariance becomes a covariance matrix and $\rho_t$ becomes a vector --- is the obvious match for wet-lab campaigns optimising stability, activity, and solubility jointly. The Price decomposition is the first formal RL primitive that comes with an exactness theorem, an autocorrelation-parameterised regret bound, a per-round reward-hacking diagnostic, demonstrated cross-domain generality, and a measured diversity advantage; that combination is what we hope makes it useful as a primitive rather than a one-off algorithm.
